10.3 The Differential of a Mapping 71

Definition 3 Relative to the original mapping (10.43) the mapping $\varphi_i : U_i \rightarrow Y$ is
called the partial mapping with respect to the variable $x_i$ at $a \in X$.

Definition 4 If the mapping (10.44) is differentiable at $x_i = a_i$, its derivative at that
point is called the partial derivative or partial differential of $f$ at $a$ with respect to
the variable $x_i$.

We usually denote this partial derivative by one of the symbols
$$
\partial_i f(a), \quad D_i f(a), \quad \frac{\partial f}{\partial x_i}(a), \quad f'_i(a).
$$

In accordance with these definitions $D_i f(a) \in \mathcal{L}(X_i; Y)$. More precisely,
$$
D_i f(a) \in \mathcal{L}(TX_{i,a_i}; TY_{f(a)}).
$$

The differential $df(a)$ of the mapping (10.43) at the point $a$ (if $f$ is differentiable at that point) is often called the total differential in this situation in order to
distinguish it from the partial differentials with respect to the individual variables.

We have already encountered all these concepts in the case of real-valued functions of $m$ real variables, so that we shall not give a detailed discussion of them. We
remark only that by repeating our earlier reasoning, taking account of Example 9 in
Sect. 9.2, one can prove easily that the following proposition holds in general.

Proposition 3 If the mapping (10.43) is differentiable at the point $a = (a_1, \dots, a_m)$
$\in X_1 \times \cdots \times X_m = X$, it has partial derivatives with respect to each variable at
that point, and the total differential and the partial differentials are related by the
equation
$$
df(a)h = d_1 f(a)h_1 + \cdots + d_m f(a)h_m, \tag{10.45}
$$
where $h = (h_1, \dots, h_m) \in TX_1(a_1) \times \cdots \times TX_m(a_m) = TX(a)$.

We have already shown by the example of numerical functions that the existence
of partial derivatives does not in general guarantee the differentiability of the function (10.43).

10.3.5 Problems and Exercises

1. a) Let $A \in \mathcal{L}(X; X)$ be a nilpotent operator, that is, there exists $k \in \mathbb{N}$ such
that $A^k = 0$. Show that the operator $(E - A)$ has an inverse in this case and that
$(E - A)^{-1} = E + A + \cdots + A^{k-1}.$

b) Let $D : \mathcal{P}[x] \rightarrow \mathcal{P}[x]$ be the operator of differentiation on the vector space
$\mathcal{P}[x]$ of polynomials. Remarking that $D$ is a nilpotent operator, write the operator
$\exp(aD)$, where $a \in \mathbb{R}$, and show that $\exp(aD)(P(x)) = P(x + a) =: T_a(P(x)).$